\documentclass[11pt,a4paper]{article}

\usepackage[margin=2.5cm]{geometry}
\usepackage[T1]{fontenc}
\usepackage{graphicx}
\usepackage{amsmath}
\usepackage{booktabs}
\usepackage[hidelinks]{hyperref}
\graphicspath{{figures/}}

\title{Advanced Visual AI Coursework: Image Super-Resolution}
\author{Edward Boileau (2209465)}
\date{}

\begin{document}
\maketitle

\begin{abstract}
% Summarise the problem, methods, experiments, main findings and conclusions.
% Write this once the experiments are complete.
\noindent\textit{Abstract to be written.}

\noindent Code: \url{https://github.com/edwardboileau/cw_advai_26}
\end{abstract}

\section{Introduction}
% Introduce image super-resolution, its challenges and the aims of this study.

\section{Selected Methods}
% Justify the choice of methods.

\subsection{Method 1}
% Replace the heading with the method name. Explain the method and cite sources.

\subsection{Method 2}
% Explain the second method and cite sources.

% Uncomment if completing the optional challenge.
% \subsection{Method 3}
% Explain the method identified from literature outside the taught curriculum.

\section{Experimental Design}

\subsection{Dataset and Metrics of Evaluation}
% Describe DIV2K, the selected image splits, preprocessing and image pairing.
% Define PSNR, SSIM and LPIPS, including ranges and evaluation conventions.

\subsection{Training}
% Describe model settings, objectives, optimisers, stopping rules, seeds,
% hardware and runtime. Distinguish dataset training from per-image fitting.

\subsection{Validation}
% Explain how settings are chosen and how evaluation images are kept separate.
% Do not optimise reconstruction parameters using evaluation HR targets.

\subsection{Qualitative and Quantitative Comparison of Methods}

\subsubsection{Bicubic $\times 8$ Super-Resolution}
% Compare the methods and a bicubic baseline using metrics and image crops.

\subsubsection{Impact of Noise Degradation}
% Specify where noise is added, its distribution, levels, units and clipping.
% Compare low and high noise levels using the same evaluation protocol.

\subsubsection{Impact of Downscaling}
% Further downscale the supplied x8 inputs by 2 and reconstruct at x16.
% Explain changes to methods and compare performance with the x8 experiments.

\section{Conclusions}
% Summarise findings, limitations and possible future work.

% Add BibTeX entries to references.bib and cite them using \cite{citation-key}.
% Uncomment the following lines after adding your first citation.
% \bibliographystyle{plain}
% \bibliography{references}

\end{document}
